\section{Problem 4: Rolling Coin on a Flat Surface}
\begin{figure} [ht]
\centering
\includegraphics[width=0.7\textwidth]{figures/Problem_4.png}
\caption{Rolling coin system.}
\end{figure}
We use the same generalized coordinates and reference frames from HW01:
\[
q = [q_1, q_2, q_3, q_4, q_5]^T
\]
where:
\begin{itemize}
    \item $q_1, q_5$ are the planar coordinates of the contact point,
    \item $q_2$ is the lean angle,
    \item $q_3$ is the spin angle,
    \item $q_4$ is the heading angle.
\end{itemize}

The angular velocity of the coin expressed in frame $B$ is:
\begin{align}
\bm{\omega}_{B/N}
&=
\dot q_2 \cos q_3 \, \mathbf{b}_1
\\
&\quad
+
(\dot q_4 \sin q_2 - \dot q_2 \sin q_3)\mathbf{b}_2
\\
&\quad
+
(\dot q_3 + \dot q_4 \cos q_2)\mathbf{b}_3
\end{align}

The standard inertia tensor for a thin disk is:
\[
I_1 = I_2 = \frac14 mR^2,
\qquad
I_3 = \frac12 mR^2
\]

\subsection*{(a) Total Kinetic Energy}

The total kinetic energy consists of:
\begin{itemize}
    \item translational kinetic energy of the center of mass,
    \item rotational kinetic energy about the center.
\end{itemize}

The translational velocity of the center is:
\[
\mathbf{v}_C
=
\dot q_1\mathbf{a}_x
+
\dot q_5\mathbf{a}_y
\]

Hence:
\[
|\mathbf{v}_C|^2
=
\dot q_1^2 + \dot q_5^2
\]

The translational kinetic energy is:
\[
T_{trans}
=
\frac12 m(\dot q_1^2+\dot q_5^2)
\]

The rotational kinetic energy is:
\[
T_{rot}
=
\frac12\bm{\omega}^T\mathbf{I}_C\bm{\omega}
\]

Substituting the angular velocity components:
\begin{align}
T_{rot}
&=
\frac12 I_1\dot q_2^2\cos^2 q_3
\\
&\quad
+
\frac12 I_1(\dot q_4\sin q_2-\dot q_2\sin q_3)^2
\\
&\quad
+
\frac12 I_3(\dot q_3+\dot q_4\cos q_2)^2
\end{align}

Therefore the total kinetic energy is:
\begin{align}
T
&=
\frac12 m(\dot q_1^2+\dot q_5^2)
\\
&\quad
+
\frac12 I_1\dot q_2^2\cos^2 q_3
\\
&\quad
+
\frac12 I_1(\dot q_4\sin q_2-\dot q_2\sin q_3)^2
\\
&\quad
+
\frac12 I_3(\dot q_3+\dot q_4\cos q_2)^2
\end{align}

\subsection*{(b) Potential Energy}

The height of the center of the coin is:
\[
z_C = R\cos q_2
\]

Hence the gravitational potential energy is:
\[
V = mgR\cos q_2
\]

\subsection*{(c) Virtual Work and Ideal Constraints}

The rolling without slipping constraints require that the contact point velocity vanish:
\[
\mathbf{v}_P = 0
\]

The contact forces consist of:
\begin{itemize}
    \item normal force $N$,
    \item friction forces $f_x, f_y$.
\end{itemize}

The virtual displacement of the contact point is:
\[
\delta\mathbf{r}_P = 0
\]
for admissible virtual displacements.

Therefore the virtual work done by the constraint forces is:
\begin{align}
\delta W_c
&=
\mathbf{F}_c\cdot\delta\mathbf{r}_P
\\
&=0
\end{align}

Hence the rolling constraints are ideal constraints.

\subsection*{(d) Lagrange Equations with Multipliers}

The rolling constraints are:
\begin{align}
\dot q_1 &= f_1(q,\dot q)
\\
\dot q_5 &= f_2(q,\dot q)
\end{align}

Define the constraint equations:
\begin{align}
\Phi_1 &= \dot q_1 - f_1(q,\dot q)=0
\\
\Phi_2 &= \dot q_5 - f_2(q,\dot q)=0
\end{align}

Introduce Lagrange multipliers $\lambda_1$ and $\lambda_2$.

The constrained Lagrange equations are:
\[
\frac{d}{dt}
\left(
\frac{\partial L}{\partial \dot q_i}
\right)
-
\frac{\partial L}{\partial q_i}
=
\lambda_1\frac{\partial\Phi_1}{\partial\dot q_i}
+
\lambda_2\frac{\partial\Phi_2}{\partial\dot q_i}
\]

for:
\[
i = 1,2,3,4,5
\]

Explicitly:

\subsubsection*{Equation for $q_1$}

\[
m\ddot q_1 = \lambda_1
\]

\subsubsection*{Equation for $q_2$}

\[
\frac{d}{dt}\left(\frac{\partial L}{\partial\dot q_2}\right)
-
\frac{\partial L}{\partial q_2}
=
-\lambda_1\frac{\partial f_1}{\partial\dot q_2}
-
\lambda_2\frac{\partial f_2}{\partial\dot q_2}
\]

\subsubsection*{Equation for $q_3$}

\[
\frac{d}{dt}\left(\frac{\partial L}{\partial\dot q_3}\right)
-
\frac{\partial L}{\partial q_3}
=
-\lambda_1\frac{\partial f_1}{\partial\dot q_3}
-
\lambda_2\frac{\partial f_2}{\partial\dot q_3}
\]

\subsubsection*{Equation for $q_4$}

\[
\frac{d}{dt}\left(\frac{\partial L}{\partial\dot q_4}\right)
-
\frac{\partial L}{\partial q_4}
=
-\lambda_1\frac{\partial f_1}{\partial\dot q_4}
-
\lambda_2\frac{\partial f_2}{\partial\dot q_4}
\]

\subsubsection*{Equation for $q_5$}

\[
m\ddot q_5 = \lambda_2
\]

The multipliers represent the unknown rolling constraint forces.

\subsection*{(e) Reduced Equations Using Independent Speeds}

The rolling constraints allow elimination of:
\[
\dot q_1,
\qquad
\dot q_5
\]

Thus the independent generalized speeds are:
\[
\dot q_2,
\qquad
\dot q_3,
\qquad
\dot q_4
\]

Substituting the constraint equations into the kinetic energy yields a reduced Lagrangian:
\[
L_r(q_2,q_3,q_4,\dot q_2,\dot q_3,\dot q_4)
\]

The reduced equations of motion become:
\[
\frac{d}{dt}
\left(
\frac{\partial L_r}{\partial\dot q_i}
\right)
-
\frac{\partial L_r}{\partial q_i}
=0
\]

for:
\[
i=2,3,4
\]

Compared with the multiplier formulation:
\begin{itemize}
    \item the reduced equations contain fewer variables,
    \item the constraint forces do not appear explicitly,
    \item the equations are expressed only in terms of independent generalized speeds.
\end{itemize}

\subsection*{(f) Steady Circular Rolling}

For steady circular rolling:
\[
\dot q_2 = 0,
\qquad
\ddot q_2 = 0
\]
\[
\dot q_3 = \Omega_s = \text{constant}
\]
\[
\dot q_4 = \Omega_h = \text{constant}
\]

The center of mass moves in a circle of radius $\rho$.

The translational speed is:
\[
v = \rho\Omega_h
\]

The required centripetal acceleration is:
\[
a_c = \frac{v^2}{\rho}
\]

Balancing moments about the contact point gives:
\[
mgR\sin q_2
=
\frac{mv^2}{\rho}R\cos q_2
\]

Therefore:
\[
\boxed{
\tan q_2
=
\frac{v^2}{g\rho}
}
\]

This gives the equilibrium relation between:
\begin{itemize}
    \item lean angle,
    \item rolling speed,
    \item turning radius.
\end{itemize}

\subsection*{(g) Conservation of Mechanical Energy}

The total mechanical energy is:
\[
E = T + V
\]

Since:
\begin{itemize}
    \item gravity is conservative,
    \item rolling constraints are ideal,
    \item no slipping occurs,
    \item no frictional dissipation exists,
\end{itemize}

we have:
\[
\frac{dE}{dt}=0
\]
